\section{Supplementary Specification and Results}
\label{app:supplementary}

This appendix holds material demoted from the body for length: specification detail that
Table~\ref{tab:submodels} states in summary, the caveats that Section~\ref{ch:discussion} does not
argue at length, and the sensitivity analyses that Section~\ref{sec:results-robustness} reports
only in summary.

\subsection{Submodel 9: The Credit-Granting Gate}
\label{app:gate}

The lender applies a residual-income test and not a debt-service-to-income cap, because Regulation
23A \cite{sa_ncr_affordability_2015} prescribes no such ratio. Credit is granted only where gross
income, less statutory deductions, less the Regulation 23A necessary expenses for the applicant's
income band, less visible existing obligations, covers the new instalment. The resulting ceiling
varies with income, from 10.4\% of income at R900 per month to 83.2\% at R7{,}712 per month, and
both worked examples published with the regulation are asserted in the verification suite.

Regulation 23A applies to an individual consumer and is applied here to household income, because
\ac{NIDS} is a household survey. For a multi-earner household this makes the test more permissive
than the regulation would be in practice, so the model's lender is somewhat more willing to lend
than a compliant lender would be.

\subsection{Submodel 17: Activation Order}
\label{app:activation}

Households act in random asynchronous order, reseeded every tick. The peer channel of Submodel~11
is activation-order independent by construction, since it reads the reference-group adoption share
lagged one tick, so every household within a tick sees the same settled share regardless of when it
acts. The model is nonetheless re-run under uniform and under fully synchronous activation, and any
material movement is reported as activation-order sensitivity, since Comer and Loerch
\cite{comer2013activation} and Alizadeh and Cioffi-Revilla \cite{alizadeh2015activation} both find
emergent population behaviour to be sensitive to the activation regime.

Credit in this model is scarce, rationed by the affordability gate and by per-platform limits, so a
fixed order would hand the same households first claim on that scarce resource every tick, which is
why random reseeded activation is the baseline rather than a fixed sweep.

\subsection{Further Caveats}
\label{app:caveats}

Table~\ref{tab:further-caveats} records seven limitations that are real but do not bound the
thesis's central claims. The five that do are argued in Section~\ref{ch:discussion}.

\begin{table}[H]
\centering
\caption{Further limitations, with the claim each one bounds.}
\label{tab:further-caveats}
\begin{tabular}{L{3.2cm}L{6.4cm}L{4.4cm}}
\toprule
\textbf{Caveat} & \textbf{Statement} & \textbf{What it bounds} \\
\midrule
\ac{FinScope} vintage & No 2017 wave exists, so the 2019 wave supplies the four
financial-inclusion flags for a population otherwise fixed at 2017. The imported variables are
categorical and need no deflation & Banked status gates \ac{BNPL} eligibility, so the 83.1\%
access ceiling carries a two-year vintage gap \\
\addlinespace
Liquid savings proxy & Derived from a \ac{NIDS} field measuring reported financial assets rather
than an accessible cash balance, and winsorised to remove implausible outliers & Validation
pattern~4 measures the share reaching zero savings, so the model is judged partly on its weakest
input \\
\addlinespace
Servicing constructed & Statutory maxima are ceilings and the \ac{CCMR} publishes no realised
rates, so opening servicing is an upper bound; loan terms are assumed with a six-month floor; the
product mix implied by assigning sub-sector rates to a consolidated balance is itself an assumption
& Opening balance sheets, and therefore the level of baseline distress \\
\addlinespace
Regulation 23A unit & The regulation prescribes a minimum-expense norm for an individual consumer
and the model applies it to household income & For multi-earner households the gate is more
permissive than in practice \\
\addlinespace
Resample distortion & Weighted resampling raises the per-capita income Gini from 0.651 to 0.671,
passing its tolerance with little margin. Measured, not corrected & If anything, modestly
overstates distress among the lowest-income agents \\
\addlinespace
Weak interaction & Each household observes its group's aggregate adoption share, never another
household individually. No South African data exists on who influences whom & The social channel
is a mean-field one, as in Granovetter's threshold models, not an explicit network \\
\addlinespace
\ac{CPI} derivation & The deflator is taken from published Stats SA releases rather than
re-derived from the underlying index table, reproducing the published index to within 0.8\% &
All \ac{BNPL} Rand parameters, which are published at current vintage and deflated once to 2017
Rands \\
\bottomrule
\end{tabular}
\end{table}

\subsection{Supplementary Sensitivity Results}
\label{app:sensitivity}

% CONTENT REQUIRED, all as tables, generated from results/summary/:
% - Sobol first-order and total-order indices for the default rate (sobol_indices.json,
%   sobol_default_rate.png).
% - Activation-order re-runs: random asynchronous vs uniform vs fully synchronous.
% - Population-size stability at 1,000 / 5,000 / 10,000 agents.
% - Minimum-payer share swept 0.20-0.40, bracketing the transferred US estimate of 0.29.
% - Income-shock magnitude sweep.
% - Full platform-count detail, 1 through 6 (the body reports endpoints only).
% - The beta=0 intervention panel (rq3_interventions_beta0.png).
% - Validation patterns 3 and 4 in full: pattern3_dti_by_quintile.png and the
%   zero-savings comparison, both demoted from the body per decision D4.

\subsection{Further Extensions}
\label{app:extensions}

Two extensions were considered and are recorded here rather than in
Section~\ref{sec:future-work}, being further from the thesis's own results.

\textbf{Macroeconomic dynamics.} Submodel~16 holds the environment static to protect the
experimental control. A successor design could relax this, since \ac{BNPL} grew in South Africa
through a period of rising rates, and the interaction between \ac{BNPL} stress and a downturn is
what a regulator would want to know. The injection effect and the macro effect would then have to
be separately identified, which requires a more elaborate design than a single counterfactual
injection.

\textbf{Competition among traditional lenders.} The traditional side is a single non-adaptive stub.
Hamill et al. \cite{hamill2023creditcard} show that promotional strategies chosen by
profit-maximising lenders raise borrower indebtedness, so a competing, adapting traditional sector
would add a driver of distress that this model holds fixed.

